\section*{Introducción}

Las inestabilidades de Rayleigh-Taylor (iRT) son fenómenos que ocurren cuando un sistema con dos fluídos de diferentes densidades, sujetos a influencias gravitatorias (o bajo algún campo potencial que actúe sobre ellos),
son dispuestos de manera tal que el más pesado ($\rho_h$) se posa sobre el liviano ($\rho_l$),
volviendo el sistema inestable bajo perturbaciones, sobre todo en o cerca a la interface. \cite{Hillier}
El estudio de este fenómeno empieza en 1883 con Rayleigh, y es seguido en los 50's con Taylor. En ellos, se utilizan las ecuaciones de Euler


